\documentclass[11pt,a4paper]{report}

\usepackage{float} %dzieki temu pakietowi rysunki pojawiaja sie tam gdzie chce po uzyciu opcji H w srodowisku figure
\usepackage[font=small,labelfont=bf]{caption}
\usepackage{natbib}
\usepackage{import}
\usepackage[tight]{subfigure}
\usepackage{tikz}
\usepackage{pgfplots}
\usepackage{polski}
\usepackage[cp1250]{inputenc}
\usepackage{indentfirst}
\usepackage{color}
\usepackage{graphics}
\usepackage{geometry}
% ---- PAGE LAYOUT ----
\geometry{a4paper, top=1.5cm, bottom=1.5cm, left=1.5cm,
right=1.5cm, nohead}
\usepackage{amsmath,amssymb,amsfonts}
\begin{document}
\begin{center}
\large{\textbf{PRACA DOMOWA III}}
\end{center}

\noindent
imi� i nazwisko ........................................\\



\noindent
\textbf{Zadanie 1}\\ 
\textbf{(a)} Wyznacz pole zawarte pomi�dzy liniami $y = 3x - x^2$ oraz $y = 2x$.\\ 
\textbf{(b)} Wyznacz pole zawarte pomi�dzy liniami $y = x^2 - 4x + 5$ oraz $y = 5 - x^2$.\\ 
\textbf{(c)} Wyznacz obj�to�� obszaru powsta�ego przez obr�t linii $y = 1-x^2$ wok� osi Ox, gdy $-1\leq x\leq 1 $.\\
\textbf{(d)} Wyznacz obj�to�� obszaru powsta�ego przez obr�t linii $y = \frac{1}{x^2}$ wok� osi Ox, gdy $1\leq x\leq 2$.\\


\noindent
\textbf{odpowiedzi}:\\ 
\textbf{(a)} $S = \frac 16$, \textbf{(b)} $S = \frac 83$, \textbf{(c)} $V = \frac{16}{15}\pi$,  \textbf{(d)} $V = \frac{7}{24}\pi$.\\

\begin{figure}[H]
	\centering
	\begin{tikzpicture}
	\foreach \x in {0,...,34}
	\draw (0.5*\x,0)--(0.5*\x,18cm);
	\foreach \y in {0,...,36}
	\draw (0,0.5*\y)--(17cm,0.5*\y);
	\end{tikzpicture}
\end{figure}

\begin{figure}[H]
	\centering
	\begin{tikzpicture}
	\foreach \x in {0,...,34}
	\draw (0.5*\x,0)--(0.5*\x,27cm);
	\foreach \y in {0,...,54}
	\draw (0,0.5*\y)--(17cm,0.5*\y);
	\end{tikzpicture}
\end{figure}

\pagebreak

\noindent
\textbf{Zadanie 2} Rozwi�� uk�ady r�wna� metod� Gaussa eliminacji:\\ 
$$
\textbf{(a)}\ 
\left \{
\begin{array}{l}
	x + 2y - 2z = 2\\
	3x + 2y + z = -3\\
	2x - 2y + z = 7\\
\end{array}
\right .,\ 
\textbf{(b)}\ 
\left \{
\begin{array}{l}
	3x + y  = 1\\
	-x + 2y + z = -2\\
	2x + 3y + z = -1\\
\end{array}
\right .,\ 
\textbf{(c)}\ 
\left \{
\begin{array}{l}
	2x + y - z = 2\\
	3x - 2y + z = 3\\
	-x + 3y - 2z = 1\\
\end{array}
\right .,\  
\textbf{(d)}\ 
\left \{
\begin{array}{l}
	x + 2y + t = -3\\
	y + z + 2t = -3\\
	2x + z + t = 2\\
	x + y + 2z = 4
\end{array}
\right . 
$$


\vspace{0.25cm}
\noindent
\textbf{odpowiedzi}: \\
\textbf{(a)} $x = 2, y= -3, z = -3$, 
\textbf{(b)} $y = 1 - 3x, z =  -4 + 7x, x\in\mathbb{R}$,  
\textbf{(c)} uk�ad sprzeczny, \\
\textbf{(d)} $x = 1, y = -1, z = 2, t = -2$.\\

\begin{figure}[H]
	\centering
	\begin{tikzpicture}
	\foreach \x in {0,...,34}
	\draw (0.5*\x,0)--(0.5*\x,19cm);
	\foreach \y in {0,...,38}
	\draw (0,0.5*\y)--(17cm,0.5*\y);
	\end{tikzpicture}
\end{figure}

\begin{figure}[H]
	\centering
	\begin{tikzpicture}
	\foreach \x in {0,...,34}
	\draw (0.5*\x,0)--(0.5*\x,27cm);
	\foreach \y in {0,...,54}
	\draw (0,0.5*\y)--(17cm,0.5*\y);
	\end{tikzpicture}
\end{figure}

\pagebreak

\noindent
\textbf{Zadanie 3} Rozwi�� uk�ady r�wna� lub wyznacz wskazan� niewiadom� stosuj�c wzory Cramera:
$$
\textbf{(a)}\ 
\left \{
\begin{array}{l}
	x + 2y - 2z = 2\\
	3x + 2y + z = -3\\
	2x - 2y + z = 7\\
\end{array}
\right .,\ 
\textbf{(b)}\ 
\left \{
\begin{array}{l}
	x + 2y + t = -3\\
	y + z + 2t = -3\\
	2x + z + t = 2\\
	x + y + 2z = 4
\end{array}
\right .\ z =?,\ 
\textbf{(c)}\ 
\left \{
\begin{array}{l}
	x + y - z + t = 2\\
	2x - y - 2z + t = -2\\
	x + 2y - 2z + t = 3\\
	2x - y + z + t = 1\\
\end{array}
\right .\ y =?,\ ,\ 
$$
$$
\textbf{(d)}\ 
 \left \{
 \begin{array}{l}
 	2x + y + z + t = 1\\
 	x + 2y + z + t = 2\\
 	x + y + 2z + t = 5\\
 	x + y + z + 2t = 2
 \end{array}
 \right .\ z =?,\ 
$$

\vspace{0.25cm}
\noindent
\textbf{odpowiedzi}: 
\textbf{(a)} $x = 2, y= -3, z = -3$, 
\textbf{(b)} $z = 2$,  
\textbf{(c)} $y = 2$, 
\textbf{(d)} $z = 3$.\\ 

\begin{figure}[H]
	\centering
	\begin{tikzpicture}
	\foreach \x in {0,...,34}
	\draw (0.5*\x,0)--(0.5*\x,19cm);
	\foreach \y in {0,...,38}
	\draw (0,0.5*\y)--(17cm,0.5*\y);
	\end{tikzpicture}
\end{figure}

\begin{figure}[H]
	\centering
	\begin{tikzpicture}
	\foreach \x in {0,...,34}
	\draw (0.5*\x,0)--(0.5*\x,27cm);
	\foreach \y in {0,...,54}
	\draw (0,0.5*\y)--(17cm,0.5*\y);
	\end{tikzpicture}
\end{figure}

\end{document}